% Avsnitt 6.9, ur lectures/L06/appendix/b_exercises.md.

\section{Övningar}
\label{c6:sec:exercises}

\begin{aside}
\textbf{Så kontrollerar du ditt arbete.} Varje övning nedan som ber dig skriva en VHDL-modul har en
självkontrollerande testbänk under \file{lectures/L06/exercises/}. Skriv din modul i dess katalog
\file{exercises/<module>/}, med det entitetsnamn och den \textbf{portordning} övningen anger, och
kör den sedan med GHDL \textemdash{} se \secref{c2:sec:testbenches} för de tre kommandona.
\end{aside}

% ------------------------------------------------------------------------------------------
\exerciseset{Räknare}

\exercise{Bredd och överslagspunkt}{Resonemang}
\label{c6:ex:width}
En signal deklareras som:

\begin{vhdlcode}
signal counter: natural range 0 to 1023;
\end{vhdlcode}

\xpart{a} Hur många bitar bred blir det register som det här syntetiseras till?

\xpart{b} Vilket är det största värde \code{counter} kan hålla innan den slår över tillbaka till
\code{0} vid nästa ökning?

\xpart{c} Skriv om deklarationen för en räknare som i stället behöver räkna från \code{0} upp till
\code{63}.

Följ räknarmönstret i \secref{c6:sec:counter}.

\exercise{Spåra ett överslag för hand}{Resonemang}
\label{c6:ex:trace}
Spåra en räknare av typen \code{natural range 0 to 3} för hand (ingen VHDL behövs). Anta att
räknaren startar på \code{2}, att den ökas en gång per klockflank under fem flanker i följd, och
att det inte finns någon reset däremellan.

Notera, klockflank för klockflank: värdet som hålls före flanken, och värdet som fångas på flanken.
Förklara varför räknaren återgår till \code{0} utan någon uttrycklig ”om räkningen når sitt
maxvärde, nollställ den”-logik.

\note[Observera]Det här är det \emph{syntetiserade} beteendet, det som spelar roll för hårdvara.
\Secref{c6:sec:counter} flaggar för det simuleringsförbehåll som hör till: GHDL
intervallkontrollerar en signal av typen \code{natural range 0 to 3}, så en simulering av den här
räknaren avbryter vid överslaget i stället för att rulla över. Spåra den på papper som hårdvara,
och grip efter en \code{unsigned} i full bredd om du någonsin behöver en räknare som slår över också
i simulering.

\exercise{Räknaren för hand, och biten som tar slut}{Krets}
\label{c6:ex:counterhand}
Bygg 4-bitarsräknaren från \secref{c6:sec:counter} för hand i CircuitVerse, och använd den för att
se spillet i stället för att ta det på förtroende.

\xpart{a} Bygg den: ett 4-bitars register (fyra D-vippor som delar en klocka), en adderare med
registrets utgångar på den ena ingången och en konstant \code{1} på den andra, adderarens summa
kopplad tillbaka till registrets \code{D}-ingångar, och fyra utgångselement eller lysdioder, en per
bit.

\xpart{b} Kör den med en klockperiod på \code{1000 ms} och notera vad du ser:
\begin{itemize}
  \item Skriv ner sekvensen av räkningar över sexton klockflanker i följd, med start från
    \code{0000}.
  \item Håll ögonen på \code{carry_out} på den översta adderaren medan räkningen rullar från
    \code{1111} till \code{0000}: under hur många av de sexton flankerna är den hög, och vad är
    kopplat till den?
  \item Förklara, med en mening, vart den femte biten tar vägen, och varför registret lagrar
    \code{0000} snarare än \code{10000}.
\end{itemize}

\xpart{c} Resonera nu om det du \emph{inte} byggde: peka på den del av din ritning som avgör att
räknaren ska återgå till \code{0000}. Det finns ingen. Säg vad som gör jobbet i stället. Vad skulle
du behöva lägga till för att få den att räkna \code{0} till \code{9} och sedan slå över?
(Övning~\ref{c6:ex:counter} ber dig skriva precis det i VHDL, så spara ditt svar.)

\xpart{d} Ändra bredden till tre bitar genom att ta bort den översta vippan och smalna av
adderaren. Förutsäg den nya överslagspunkten innan du kör den, bekräfta den, och formulera regeln
som kopplar en räknares bredd till det värde den slår över vid.

\textbf{Spara den här kretsen.} \Chapref{c7:ch} bygger sin timer genom att lägga till två saker till
den, en jämförare och en nollställning, och att utgå från en räknare du redan har sett fungera är
hela poängen med övningen där.

\exercise{\code{counter}: en räknare med godtycklig radix}{VHDL}
\label{c6:ex:counter}
Skriv en entitet med namnet \code{counter} som räknar från \code{0} upp till ett steg under sin
radix, och sedan slår över tillbaka till \code{0}. Med standardvärdet på radix är det en
decimalräknare, \code{0} till \code{9}.

Entiteten har en generic:

\begin{booktable}{@{}p{4.4em}p{9.8em}p{4em}L@{}}
  \thead{Generic} & \thead{Typ} & \thead{Förval} & \thead{Betydelse} \\
  \midrule
  \code{RADIX} & \code{natural range 1 to 2**16} & \code{10} & Hur många värden räknaren går igenom
    innan den upprepar sig. Den räknar \code{0} till \code{RADIX-1} inklusive, så \code{RADIX} är
    cykelns längd, inte den största räkningen. \\
\end{booktable}

och de här portarna:

\begin{booktable}{@{}p{5.6em}p{3.4em}p{10em}L@{}}
  \thead{Port} & \thead{Riktn.} & \thead{Typ} & \thead{Betydelse} \\
  \midrule
  \code{clock} & in & \code{std_logic} & Systemklocka. \\
  \code{reset_s2_n} & in & \code{std_logic} & Aktiv låg, \textbf{redan synkroniserad}: den kommer
    från din \code{reset_sync}, aldrig rakt från en pinne. Asynkron, så den nollställer utan att
    någon klockflank är inblandad. \\
  \code{count} & out & \code{natural range 0 to RADIX-1} & Den löpande räkningen. \code{0} är ett
    värde den antar, så intervallet börjar där: reset driver den till \code{0}, och den återvänder
    dit vid varje överslag. \\
  \code{tick} & out & \code{std_logic} & En puls på en cykel vid den flank där \code{count} slår
    över från \code{RADIX-1} tillbaka till \code{0}. \\
\end{booktable}

Implementera den med en enda synkron process:
\begin{itemize}
  \item Ta med \code{clock} och \code{reset_s2_n} i känslighetslistan.
  \item Vid reset, sätt \code{count} till \code{0} och \code{tick} till \code{'0'}.
  \item Vid varje stigande klockflank: driv \code{tick} låg som utgångsläge; när \code{count} har
    nått \code{RADIX-1}, slå över \code{count} till \code{0} och pulsa \code{tick} hög; öka annars
    \code{count}.
\end{itemize}

Skriv jämförelsen mot \code{RADIX-1}, aldrig mot \code{9}. Genericen är hela skillnaden mellan en
modul och en modul du kan använda två gånger: \code{RADIX} är fastlagd innan syntesen körs, så
jämföraren den bygger blir inte större än en hårdkodad, och räknarens register dimensioneras efter
den radix du ber om. Det är samma argument som \secref{c4:sec:generics} för fram för
\code{button_sync}:s \code{COUNT}, och skälet till att \chapref{c7:ch}:s \code{timer} tar sitt
målvärde som en generic snarare än som en konstant.

Till skillnad från det gratis överslaget hos en räknare av typen \code{natural range 0 to 15}
behöver en räknare med godtycklig radix en uttrycklig jämförelse, eftersom \code{RADIX-1} i allmänhet
inte är en tvåpotensgräns. Lägg märke till vad som händer när den är det: med \code{RADIX = 16}
skrivs jämförelsen fortfarande ut, men hårdvaran under behöver den inte längre.

\note[Tips]Processen behöver läsa den löpande räkningen (för att jämföra den och för att öka den),
så håll räkningen i en intern \code{signal count_s: natural range 0 to RADIX-1;}, kör logiken på
\code{count_s}, och driv utgången med \code{count <= count_s;} utanför processen, precis som
\file{d_flip_flop.vhd} använder \code{q_s} i \chapref{c3:ch}. Det är inget stilval: varje
\code{ghdl}-kommando i den här boken skickar \code{--std=93}, och under VHDL-93 kompilerar det inte
alls att läsa en \code{out}-port. Senare standarder mjukar upp det, vilket är varför du kan ha sett
det göras på annat håll.

Följ räknarmönstret i \secref{c6:sec:counter}, och den driv-låg-och-pulsa-sedan-form som
\code{data_ready} har i \secref{c6:sec:rx8}.

\modulefig[0.68]{lectures/L06/appendix/images/counter.png}

\begin{selfcheck}
\textbf{Självkontroll.} Döp din entitet till \code{counter}, med genericen \code{RADIX}
(\code{natural range 1 to 2**16}, standardvärde \code{10}), ingångarna \code{clock},
\code{reset_s2_n} och utgångarna \code{count} (\code{natural range 0 to RADIX-1}) och \code{tick}
(\code{std_logic}), deklarerade i den ordningen; dess testbänk finns i \file{exercises/counter/}.
Den elaborerar \textbf{två} instanser, en med \code{RADIX = 10} och en med \code{RADIX = 4}, och kör
vardera genom två fulla cykler, och kontrollerar att \code{tick} är en puls på en cykel snarare än
en nivå och att den aldrig utlöser halvvägs igenom. En arkitektur som tyst hårdkodar \code{10}
klarar den första instansen och fallerar på den andra, vilket är poängen med att kontrollera två.
\end{selfcheck}

% ------------------------------------------------------------------------------------------
\exerciseset{Skiftregister}

\exercise{Fyra bitar in i ett SIPO-register}{Resonemang}
\label{c6:ex:sipotrace}
Ett 8-bitars SIPO-skiftregister startar på \code{00000000} och använder idiomet från
\secref{c6:sec:sipo}:

\begin{vhdlcode}
shift_reg <= shift_reg(6 downto 0) & serial_in;
\end{vhdlcode}

Det tar emot den seriella bitströmmen \code{1, 0, 1, 1} (en bit per klockflank, i den ordningen)
över fyra stigande klockflanker i följd.

\xpart{a} Skriv ut hela det 8-bitars innehållet i \code{shift_reg} efter var och en av de fyra
klockflankerna.

\xpart{b} Efter de fyra flankerna: vilka fyra bitar i \code{shift_reg} håller de mottagna data, och
i vilken ordning ligger de, i förhållande till den ordning bitarna anlände i?

\exercise{Fyra bitar ut ur ett PISO-register}{Resonemang}
\label{c6:ex:pisotrace}
Ett 4-bitars PISO-register, som följer idiomet från \secref{c6:sec:piso}, är parallellt laddat med
\code{parallel_in = "1010"} (\code{load = '1'} under en klockflank), sedan skiftat med
\code{shift = '1'} under de följande fyra klockflankerna, med \code{serial_in} hållen på \code{'0'}
hela tiden.

Kom ihåg hållet från \secref{c6:sec:shift}: bitarna flyttar mot MSB, \code{serial_out} är bit 3
(värdet som är på väg att lämna), och \code{serial_in} kommer in vid bit 0.

\xpart{a} Skriv ut registrets innehåll och \code{serial_out} vid varje steg, en rad per klockflank:
raden för laddningsflanken, före all skiftning, och sedan en rad för var och en av de fyra
skiftflankerna.

\xpart{b} I vilken ordning dyker de fyra laddade bitarna upp på \code{serial_out}, i förhållande
till sina positioner i \code{parallel_in}?

\xpart{c} Vad är registrets slutliga innehåll efter alla fyra skiftningarna, och varför?

\note[Tips]Du behöver ingen riktig hårdvara för att kontrollera Övning~\ref{c6:ex:sipotrace} och
\ref{c6:ex:pisotrace} för hand. Gå igenom dem precis som du skulle spåra en sanningstabell: en
klockflank, en rad, i taget.

\exercise{\code{sipo8} och \code{piso8}}{VHDL}
\label{c6:ex:shiftregs}
Implementera kapitlets två användbara skiftregisterutföranden i VHDL.

\xpart{a}\parttitle{\code{sipo8}.} Skriv en entitet med namnet \code{sipo8}, ett register med serie
in och parallell ut.

\begin{booktable}{@{}p{6.6em}p{3.4em}p{11em}L@{}}
  \thead{Port} & \thead{Riktn.} & \thead{Typ} & \thead{Betydelse} \\
  \midrule
  \code{clock} & in & \code{std_logic} & Systemklocka. \\
  \code{reset_s2_n} & in & \code{std_logic} & Aktiv låg, \textbf{redan synkroniserad}. Asynkron.
    Nollställer hela registret, så \code{parallel_out} läser \code{"00000000"} medan den hålls
    låg. \\
  \code{serial_in} & in & \code{std_logic} & Den inkommande biten. \\
  \code{parallel_out} & out & \code{std_logic_vector(7 downto 0)} & Registrets innehåll. \\
\end{booktable}

Följ SIPO-idiomet från \secref{c6:sec:sipo}.

\modulefig[0.68]{lectures/L06/appendix/images/sipo8.png}

\xpart{b}\parttitle{\code{piso8}.} Skriv en entitet med namnet \code{piso8}, ett register med
parallell in och serie ut.

\begin{booktable}{@{}p{6.6em}p{3.4em}p{11em}L@{}}
  \thead{Port} & \thead{Riktn.} & \thead{Typ} & \thead{Betydelse} \\
  \midrule
  \code{clock} & in & \code{std_logic} & Systemklocka. \\
  \code{reset_s2_n} & in & \code{std_logic} & Aktiv låg, \textbf{redan synkroniserad}. Asynkron.
    Nollställer hela registret, så \code{serial_out} läser \code{'0'} medan den hålls låg. \\
  \code{load} & in & \code{std_logic} & Fånga \code{parallel_in} vid nästa stigande flank. \\
  \code{shift} & in & \code{std_logic} & Skifta en position vid nästa stigande flank. \\
  \code{serial_in} & in & \code{std_logic} & Biten som matas in vid bit 0 vid en skiftning. \\
  \code{parallel_in} & in & \code{std_logic_vector(7 downto 0)} & Värdet som fångas vid en
    laddning. \\
  \code{serial_out} & out & \code{std_logic} & Registrets översta bit. \\
\end{booktable}

Följ PISO-idiomet från \secref{c6:sec:piso}: vid en laddningsflank (\code{load = '1'}), fånga
\code{parallel_in}; vid en skiftflank (\code{shift = '1'}), flytta varje bit en position mot MSB
och mata in \code{serial_in} vid bit 0; driv \code{serial_out} från bit 7, så att en laddad byte
lämnar med den mest signifikanta biten först; och låt \code{load} ha företräde framför \code{shift}
när båda är höga.

\modulefig[0.68]{lectures/L06/appendix/images/piso8.png}

\xpart{c} Namnge, för vart och ett av de fyra skiftregisterutförandena (SISO, SIPO, PISO, PIPO), en
uppgift som det är det naturliga valet för. Följ sammanfattningstabellen i \secref{c6:sec:shift}.

Härled varje idiom på nytt ur kapitlet snarare än att kopiera \code{serial_rx8}:s skiftrad.

\begin{selfcheck}
\textbf{Självkontroll.} \code{sipo8}: portarna \code{clock}, \code{reset_s2_n}, \code{serial_in}
(in) och \code{parallel_out} (\code{std_logic_vector(7 downto 0)}, out), deklarerade i den
ordningen; testbänk i \file{exercises/sipo8/}. \code{piso8}: portarna \code{clock},
\code{reset_s2_n}, \code{load}, \code{shift}, \code{serial_in}, \code{parallel_in}
(\code{std_logic_vector(7 downto 0)}) och \code{serial_out} (out), deklarerade i den ordningen;
testbänk i \file{exercises/piso8/}.

\textbf{Spara \code{piso8}.} \Chapref{c8:ch}:s capstone med sändaren instansierar den inte, men den
är värd att läsa om där för idiomets skull, eftersom den maskinen skiftar ut byten inline.
\end{selfcheck}

% ------------------------------------------------------------------------------------------
\exerciseset{Att läsa seriemottagaren}

Till det genomarbetade exemplet hör en referenstestbänk. Kör den från exemplets egen katalog innan
du börjar:

\begin{shell}
cd lectures/L06/serial_rx8
ghdl -a --std=93 serial_rx8.vhd serial_rx8_tb.vhd
ghdl -e --std=93 serial_rx8_tb
ghdl -r --std=93 serial_rx8_tb --assert-level=error --stop-time=10ms
\end{shell}

\exercise{Mottagarens beteende, och tre riktiga buggar}{Simulering}
\label{c6:ex:rx8}
Redogör för mottagarens beteende, med hänvisning till \secref{c6:sec:rx8}.

\xpart{a} Bitströmmen \code{1, 0, 1, 1, 0, 0, 1, 0} anländer, en bit per aktiverad klockflank. Vad
ligger på \code{data_out} under den cykel där \code{data_ready} pulsar? Vilken anländande bit hamnar
i \code{data_out(7)}, och vilken i \code{data_out(0)}?

\xpart{b} \code{shift_enable} går låg under tre klockcykler efter den fjärde biten, och sedan hög
igen under de återstående fyra. Producerar mottagaren fortfarande samma byte? Förklara vad
biträknaren gör under uppehållet. Vad skulle i stället gå fel om \code{shift_enable} grindade bara
skiftregistret och inte räknaren?

\xpart{c} Gör var och en av följande ändringar i din egen kopia av \file{serial_rx8.vhd}, en i
taget, och \textbf{förutsäg} vad referenstestbänken kommer att rapportera innan du kör den. Kör den
sedan och redogör för resultatet:
\begin{itemize}
  \item Ändra skiftraden till det spegelvända hållet,
    \code{shift_reg <= serial_in & shift_reg(7 downto 1);}.
  \item Ta bort det villkorslösa \code{data_ready <= '0';} högst upp i den klockade grenen.
  \item Ändra biträknarens jämförelse från \code{7} till \code{6}.
\end{itemize}

Var och en av dem är en verklig bugg som någon har skickat ut skarpt. Säg för var och en vilken av
testbänkens kontroller som fångar den, och hur felet hade sett ut på en verklig seriell länk om
ingenting hade fångat det.
